% Generated draft: requires mathematical and source review.
\documentclass[11pt]{article}
\usepackage[margin=25mm]{geometry}
\usepackage{fontspec}
\setmainfont{Latin Modern Roman}
\usepackage{amsmath,amssymb,mathtools,unicode-math}
\setmathfont{Latin Modern Math}
\usepackage{enumitem,needspace,xurl,hyperref}
\setlength{\parindent}{0pt}
\setlength{\parskip}{5pt}
\setlength{\emergencystretch}{3em}
\newcommand{\N}{\mathbb{N}}
\newcommand{\Z}{\mathbb{Z}}
\newcommand{\Q}{\mathbb{Q}}
\newcommand{\R}{\mathbb{R}}
\newcommand{\C}{\mathbb{C}}
\newcommand{\F}{\mathbb{F}}
\newcommand{\A}{\mathbb{A}}
\newcommand{\PP}{\mathbb{P}}
\newcommand{\E}{\mathbb{E}}
\newcommand{\eps}{\varepsilon}
\newcommand{\dd}{\,\mathrm{d}}
\begin{document}
\section*{UnsolvedMath ID: 800760}
\textbf{Source problem code:} catalogue entry
\subsection{Definitions and mathematical statement}
Let \(K\) be an infinite field, with no restriction on its characteristic. Choose a separable closure \(K^{\mathrm{sep}}\), and equip its absolute Galois group
\[
G_K=\operatorname{Gal}(K^{\mathrm{sep}}/K)
\]
with the Krull topology. A profinite group is \emph{prosolvable} if every finite continuous quotient is a solvable group. Thus \(G_K\) is prosolvable precisely when \(\operatorname{Gal}(L/K)\) is solvable for every finite Galois extension \(L/K\).

A field \(K\) is \emph{ample} if, for every smooth geometrically integral curve \(C\) over \(K\), the existence of a \(K\)-rational point implies that \(C(K)\) is infinite. Here a curve is a separated scheme of finite type and dimension one over \(K\), and geometrically integral means that its base change to an algebraic closure of \(K\) is integral.

\textbf{Question.} Is it true that, for every infinite field \(K\),
\[
G_K\text{ prosolvable}
\quad\Longrightarrow\quad
\bigl(\forall\text{ smooth geometrically integral curves }C/K,\quad C(K)\ne\varnothing\Longrightarrow \#C(K)=\infty\bigr)?
\]

\textbf{Supplied background and qualifications.} The source attributes to Colliot-Th\'el\`ene the observation that a field \(K\) is ample whenever \(G_K\) is pro-\(p\) for some prime \(p\), citing the reference labels \texttt{ColliotThelene} and \texttt{JardenAmple}. A profinite group is pro-\(p\) if every finite continuous quotient is a \(p\)-group. The source notes that this covers the archimedean local fields \(\R\) and \(\C\), that pro-\(p\) groups are prosolvable, and that all local fields have prosolvable absolute Galois groups. It presents the question as a generalization of that observation. These background claims and the registry's open status are supplied source information, not an independently verified literature assessment. The bibliographic identities behind the two reference labels are not supplied.
\subsection{Short English statement}
Must every infinite field whose finite Galois extensions all have solvable Galois groups be ample? Explicitly, must every smooth geometrically integral curve over such a field have infinitely many rational points whenever it has one? The supplied registry labels this question open; this draft does not independently verify that status.
\subsection{Sources}
\begin{itemize}
\item Local UnsolvedMath registry, permanent ID \(800760\), source problem code \texttt{catalogue entry}, title: \emph{Open Problems in the Theory of Ample Fields --- Question 3}. The record attributes the question to Bary-Soroker and Fehm, with proposed year 2011, and supplies the status \emph{open}.
\item Bary-Soroker and Fehm, \emph{Open Problems in the Theory of Ample Fields} (2011), Question 3, \url{https://arxiv.org/abs/1106.1310}. The supplied provenance records access on 2026-10-16 and extraction by \texttt{arxiv\_latex}; these are supplied provenance notes, not a claim that this draft independently accessed or verified the source.
\item The supplied statement attributes the pro-\(p\) observation to Colliot-Th\'el\`ene and cites the labels \texttt{ColliotThelene} and \texttt{JardenAmple}. Full bibliographic entries for these labels were not supplied, so their identities and the precise scope of the cited result require source review.
\end{itemize}
\end{document}
